%------------------------------------------------------------------------------%
\begin{question}
  Considerando a equação
  \[
    \frac{(x+2)^2}{16} - \frac{(y-1)^2}{9} = 1
  \]
  \vspace{-\baselineskip}
  \begin{enumerate}
    \item identifique a curva
    \item determine seus principais valores e pontos
    \item esboce a curva
  \end{enumerate}
\end{question}

%------------------------------------------------------------------------------%
\begin{resolution}{Identificação}
\begin{columns}[t]
\column{0.4\textwidth}
  \onslide<+->{A equação
  \[
    \frac{(x+2)^2}{16} - \frac{(y-1)^2}{9} = 1
  \]}%
  \onslide<+->{está na forma de uma hipérbole transladada
  \[
   \frac{(x-h)^2}{a^2} - \frac{(y-k)^2}{b^2} = 1
  \]}
  \onslide<+->{Identificamos o centro}
  \[
    \onslide<+->{C = (h,  k)}
    \onslide<+->{  = (-2, 1)}
  \]

\column{0.4\textwidth}
  \onslide<+->{Os semieixos}
  \begin{align*}
    \onslide<+->{a^2 & = 16}
    \qquad
    \onslide<+->{a = 4}
    \\
    \onslide<+->{b^2 & = 9}
    \qquad
    \onslide<+->{b = 3}
  \end{align*}

  \onslide<+->{O coeficiente de $(x-h)^2$ é positivo}
  \onslide<+->{os vértices estão em uma linha horizontal}
  \begin{align*}
    \onslide<+->{V_1
    & = C - (a, 0)}
    \onslide<+->{= (-2 -4, 1)}
    \onslide<+->{= (-6, 1)}
    \\
    \onslide<+->{V_2
    & = C + (a, 0)}
    \onslide<+->{= (-2 + 4, 1)}
    \onslide<+->{= (2, 1)}
  \end{align*}
\end{columns}
\end{resolution}

%------------------------------------------------------------------------------%
\begin{resolution}{Esboço}
  \centering
  \includegraphics{I-07-Identificacao-ex-8}
\end{resolution}

%------------------------------------------------------------------------------%
